%Targeted Conference 1: 22nd EAI International Conference on Collaborative Computing: Networking, Applications and Worksharing (CCF C)
%https://collaboratecom.eai-conferences.org/2026/
%Deadline: May 1st, 2026. Regular papers should be 16-20 pages

%Targeted Conference 2: UbiComp / ISWC 2026 (CCF A) [Ryan's paper]
%https://www.ubicomp.org/ubicomp-iswc-2026/iswc-cfp/
%Deadline: May 17, 2026. ACM 6 pages

\documentclass[sigconf,natbib=false]{acmart}
\AtBeginDocument{%
  \providecommand\BibTeX{{%
    Bib\TeX}}}
\usepackage{comment}
\usepackage{subfloat}
\usepackage{subcaption}

\makeatother
%% Bibliography style
\RequirePackage[
  datamodel=acmdatamodel,
  style=acmnumeric,
  sorting=none
  ]{biblatex}

%% Declare bibliography sources (one \addbibresource command per source)
\addbibresource{reference}
%\renewcommand\footnotetextcopyrightpermission[1]{}

\copyrightyear{2026}
\acmYear{2026}
%\setcopyright{cc}
%\setcctype{by}
\acmConference[ISWC'26]{The 30th Annual International Symposium on Wearable Computers}{October 11-15, 2026}{Shanghai, China}
\acmBooktitle{The 30th Annual International Symposium on Wearable Computers, October 11-15, 2026, Shanghai, China}
%\acmDOI{10.1145/3722570.3726881}
%\acmISBN{/2025/05}


\begin{document}

%%
%% The "title" command has an optional parameter,
%% allowing the author to define a "short title" to be used in page headers.
\title{Structured In-Context Correction in MCP-Mediated Human Authority Injection for LLM-Driven Multi-Robot Rescue}

\author{Ryan~Tran, Simon~Zhang, Wenlu~Zhang, Hailu~Xu}
\affiliation{
  \institution{California State University, Long Beach}
  \city{Long Beach}
  \state{CA}
  \country{USA}
}
%\email{{ryan.tran, simon.zhang01}@student.csulb.edu}
\email{{wenlu.zhang, hailu.xu}@csulb.edu}

\author{Lingfeng Tao}
\affiliation{
  \institution{Kennesaw State University}
  \city{Atlanta}
  \state{GA}
  \country{USA}
}
\email{ltao2@kennesaw.edu}

\author{Shengjie Xu}
\affiliation{
  \institution{University of Arizona}
  \city{Tucson}
  \state{AZ}
  \country{USA}
}
\email{sjxu@arizona.edu}

\author{Yu Yin, Shuai Xu}
\affiliation{
  \institution{Case Western Reserve University}
  \city{Cleverland}
  \state{OH}
  \country{USA}
}
\email{{yu.yin, shuai.xu2}@case.edu}

\thanks{This work is partially supported by US NSF Award \#2426470.}
\thanks{Corresponding E-email: hailu.xu@csulb.edu (H.X.)}



\renewcommand{\shortauthors}{} % Optional: Remove short author line if not needed
%% By default, the full list of authors will be used in the page
%% headers. Often, this list is too long, and will overlap
%% other information printed in the page headers. This command allows
%% the author to define a more concise list
%% of authors' names for this purpose.
\renewcommand{\shortauthors}{Xu et al.}


\begin{abstract}
Large language models (LLMs) coordinated through the Model Context Protocol (MCP) offer a principled foundation for multi-robot collaboration systems, yet existing human-in-the-loop (HITL) frameworks treat human intervention as an ad-hoc, out-of-band interrupt: the operator issues a natural-language correction and the LLM replans from scratch, retaining no structured memory of the reason for that correction. Consequently, the LLM tends to repeat the same class of error once the immediate context shifts. This paper introduces Structured Authority Injection (SAI), a mechanism that promotes human overrides to first-class MCP tool calls, typed according to three semantically distinct correction categories—constraint, priority, and strategy—and propagated as durable context entries that persist across subsequent planning cycles. We implement SAI on top of an existing MuJoCo-based multi-robot rescue simulation driven by GPT-4o and evaluate it against an unstructured baseline in which overrides are appended as free-text prompts. Across 300 controlled simulation runs (3 override types × 2 injection modes × 10 scenarios × 5 repeats), SAI reduces the post-override recurrence rate of the corrected error class by 41\% and improves mission success rate by 8.3 percentage points in complex scenarios, while requiring 27\% fewer subsequent human interventions. These results demonstrate that protocol-level structuring of human authority—rather than ad-hoc textual correction—is a high-leverage design choice for safe and efficient human–robot collaboration.

\end{abstract}
%

%%
%% Keywords. The author(s) should pick words that accurately describe
%% the work being presented. Separate the keywords with commas.
\keywords{Human-in-the-loop, Multi-robot Collaboration, Model Context Protocol, Human-robot Collaboration.}

\maketitle

\section{Introduction}

Search-and-rescue operations are among the most demanding environments for autonomous robot teams. Time pressure, partial observability, dynamic hazards, and the life-critical consequences of assignment errors create conditions under which fully autonomous systems cannot yet be trusted without human oversight. At the same time, the cognitive load of supervising a fleet of robots executing many concurrent rescue jobs quickly overwhelms a single human operator. Large language models offer an appealing middle ground: by reasoning over natural-language task descriptions and issuing tool calls to robot-control APIs, an LLM can handle routine assignment decisions autonomously, escalating to the human only when it is uncertain or when constraints are violated.

The Model Context Protocol (MCP) has recently emerged as a standardised JSON-RPC 2.0 interface for connecting LLMs to external tools and data sources. Applied to robotics, MCP enables a language model to query robot status, assign jobs, and dispatch navigation commands through a uniform, schema-validated API—an architecture we have already demonstrated in a MuJoCo-based multi-robot rescue system. However, a critical gap remains: when a human operator needs to correct an LLM decision, existing systems—including our own prior work and recent frameworks such as HMCF—treat the override as an out-of-band natural-language message that is appended to the conversation history. The LLM reads the correction, replans, and moves on. Crucially, the system retains no structured record of why the correction was made, what class of error it addressed, or how future decisions should be constrained to prevent recurrence.

This paper makes the observation that human overrides in a multi-robot rescue context fall into three semantically distinct classes: (i) constraint overrides, which correct violations of physical or operational limits (e.g., battery, payload, hazard zone); (ii) priority overrides, which adjust the urgency ordering of jobs based on information the LLM did not have; and (iii) strategy overrides, which introduce spatial or coordination heuristics that should govern planning over a longer horizon. We argue that these three classes carry different information, should be represented differently in the LLM's context, and should have different persistence lifetimes. Encoding all three as free text treats them identically and discards their semantic structure.

We propose Structured Authority Injection (SAI), a lightweight extension to our MCP server that exposes three typed override tools—one per correction class. When the human operator selects an override type and fills in a structured form in the supervision interface, the corresponding MCP tool is called, and the resulting structured record is appended to a persistent authority context that is prepended to every subsequent LLM prompt for the duration of the mission. This design has two key properties: the override is a first-class protocol event (logged, schema-validated, and auditable), and its semantic content is preserved in a form the LLM can reliably attend to across many planning cycles.

We evaluate SAI against an unstructured baseline—identical architecture but overrides expressed as free-text prompt appends—in a controlled experiment using our MuJoCo rescue simulation with GPT-4o. The experiment is designed to be executable within three weeks: 10 fixed scenarios across 3 override types and 2 injection modes, each repeated 5 times, for 300 total simulation runs. Our primary metric is the Recurrence Rate (RR): the fraction of post-override decisions that repeat the same class of error as the one that triggered the override. Secondary metrics are mission success rate, override count per mission, and correction latency.

The contributions of this paper are as follows:

\begin{enumerate}
\item We introduce Structured Authority Injection (SAI), the first mechanism to promote human overrides in an LLM-driven multi-robot system to typed, schema-validated MCP tool calls, with persistent propagation across planning cycles.
\item We define a three-class taxonomy of rescue-domain overrides (constraint, priority, strategy) and design MCP tool schemas for each, together with a persistence model that governs how each class modifies the LLM's planning context.
\item We provide controlled empirical evidence that structured override injection significantly reduces error recurrence and reduces the burden on the human operator, compared to unstructured free-text correction.
\item We discuss the implications for human-robot collaboration system design, in particular the role of protocol-level structuring as a lever for improving both safety and autonomy.

\end{enumerate}

The remainder of the paper is organised as follows. Section 2 reviews related work. Section 3 describes the system architecture. Section 4 presents the SAI design. Section 5 details the experimental methodology. Section 6 reports results. Section 7 discusses implications and concludes.
\section{Related Work}

\subsection{LLM-Driven Multi-Robot Collaboration}

Several recent systems have demonstrated that LLMs can serve as high-level planners for multi-robot task allocation. COHERENT [cite] decomposes rescue missions into sub-tasks and allocates them to heterogeneous robots through a propose–execute–feedback loop, achieving strong performance on simulated disaster scenarios. RoCo [cite] and ReAd [cite] use multi-agent dialogue between LLM instances to reach consensus on robot assignments. MetaGPT [cite] formalises multi-agent LLM workflows using structured role assignments. None of these systems incorporate a human operator in the planning loop, and none use a standardised protocol layer such as MCP for tool interaction.

\subsection{Human-in-the-Loop Robot Systems}
Human oversight of robot teams has a long history in the shared-autonomy literature. Classical approaches use mixed-initiative control [cite] or sliding autonomy [cite], where the system's autonomy level is adjusted based on task difficulty or robot confidence. More recent work applies this idea to LLM-based systems: HMCF [cite] integrates GPT-4o with a human supervision interface for multi-robot rescue, allowing operators to issue natural-language corrections when the LLM's plan is unsatisfactory. Sirius [cite] provides a shared-control framework in which human corrections are used to adjust an autonomous agent's future behaviour. BrainBody-LLM [cite] uses closed-loop execution feedback—including human-in-the-loop signals—to trigger replanning in a single-robot manipulation task.

A common thread across all of these systems is that human corrections are communicated as natural language or low-level control signals, and are not structured at the protocol level. The LLM reads the correction and replans, but the system does not classify the correction semantically, does not validate it against a schema, and does not guarantee that the correction's intent will persist across subsequent planning steps. Our work addresses this gap directly.

\subsection{MCP in Robotic Systems}
MCP was introduced by Anthropic in late 2024 as a universal interface for connecting LLMs to external tools. Its adoption in robotics has accelerated rapidly: Robot MCP servers [cite] wrap ROS topics and services as MCP tools, allowing LLMs to command robots through natural language. IoT-MCP [cite] demonstrated MCP-based integration of sensor streams with LLM reasoning. MCP-Zero [cite] extends MCP with proactive tool discovery, allowing LLMs to request new capabilities on demand. To our knowledge, no prior work has used MCP as a channel for structured human authority injection—existing deployments treat MCP as a one-directional LLM-to-robot interface, not as a bidirectional human–LLM–robot coordination protocol.

\subsection{Positioning of This Work}

Table 1 positions our system relative to the closest prior work across five design dimensions. The combination of HITL override, MCP protocol layer, multi-robot rescue context, typed override classification, and persistent in-context propagation is unique to this paper.


\begin{table*}[t]
\centering
\caption{Comparison with closely related systems.}
\label{tab:comparison}

\small
\setlength{\tabcolsep}{4pt}
\renewcommand{\arraystretch}{1.1}

\begin{tabular}{lccccc}
\toprule
\textbf{System} & \textbf{HITL override} & \textbf{MCP layer} & \textbf{Multi-robot rescue} & \textbf{Structured override type} & \textbf{In-context propagation} \\
\midrule
HMCF [X] & Yes & No & Yes & No & No \\
BrainBody-LLM [X] & Yes & No & No & No & No \\
COHERENT [X] & No & No & Yes & No & No \\
Sirius [X] & Yes & No & No & No & No \\
This work & Yes & Yes & Yes & Yes & Yes \\
\bottomrule
\end{tabular}
\end{table*}
\section{System Design}

The system has four layers, shown in Figure 1: (1) a MuJoCo physics simulation of the rescue arena; (2) an MCP server that exposes robot-control tools and the new SAI override tools; (3) a GPT-4o planning agent that calls tools in an agentic loop; and (4) a human supervision interface through which the operator monitors the mission and issues structured overrides.

\subsection{MuJoCo Simulation}
The rescue arena is a bounded 2-D planar world simulated at dt = 0.002 s (500 Hz) with real-time factor 1.0, ensuring that LLM inference latency has genuine operational consequences (delayed assignments translate directly into missed deadlines). The arena contains static obstacles, a configurable number of differential-drive robots, and rescue jobs that arrive stochastically. Each job has a type (medical, structural, fire, evacuation, or hazmat), a location, a priority level (1–5), and a hard deadline in simulation seconds. Robot state—pose, velocity, battery, and current assignment—is serialised into a JSON snapshot at each decision epoch and made available to the MCP server.

\subsection{MCP Server}
The MCP server exposes two categories of tools. The first category—robot-control tools—covers the core planning operations: assign_robot_to_job, query_robot_status, dispatch_move_command, report_job_complete, and get_pending_jobs. These tools were present in our prior system. The second category—the SAI override tools—is introduced in this paper and is described in detail in Section 4. All tool calls pass through a logging middleware that records timestamps, payloads, and round-trip times, providing the raw data for our latency and error-recurrence analyses.

\subsection{Human Supervision Interface}
The supervision interface renders a live top-down view of the MuJoCo arena, a feed of recent LLM tool calls and their outcomes, and an override panel. In the baseline condition (Unstructured), the override panel contains a free-text input box; the operator types a natural-language correction which is appended to the conversation history as a user-turn message. In the SAI condition (Structured), the panel presents a three-tab form—one tab per override type—with labelled fields corresponding to the MCP tool parameters. The operator selects the appropriate tab, fills in the fields, and submits; the interface calls the corresponding MCP override tool. Both conditions are identical in every other respect, including the underlying LLM, simulation, and robot-control tools.

\subsection{Structured Authority Injection}

SAI has three components: an override taxonomy that classifies human corrections, MCP tool schemas for each class, and a persistence model that governs how each class updates the LLM's planning context.
%
% ---- Bibliography ----
%
%% Print the bibliography
\printbibliography

\end{document}
